% =====================================================================================================
\section{The method, as coded}\label{sec:method}

This section describes the method exactly as the software (an R package with a C++ core) implements it in its default
configuration. The user supplies
the data, a random seed and a computing allocation; everything else is a numerical constant or a stated function of
$(n,p)$. Table~\ref{tab:params} lists the constants that matter for the results, and Appendix~\ref{app:params} gives all
of them with their code locations.

\subsection{The model and every prior}\label{sec:prior}

\paragraph{Working model.} Given a model $S$, coefficients $\theta_S\in\R^{\abs{S}}$ and a variance $\sigma^2>0$, the
method fits $y\mid S,\theta_S,\sigma^2\sim\Normal_n(X_S\theta_S,\sigma^2I_n)$ to the standardised data. Here $\theta_S$
and $\sigma^2$ are working parameters; the intercept is removed by centring. The results of Section~\ref{sec:results} are
proved under the true model~\eqref{eq:model}, not under this working model.

\paragraph{Model space and model prior.} The method considers every model with at most $K$ predictors,
\begin{equation}\label{eq:space}
  \Om:=\{S\subseteq\set{p}:\abs{S}\le K\},\qquad K:=\min\big(p,\lfloor n/2\rfloor\big),
\end{equation}
and gives a model a prior probability that depends only on its size:
\begin{equation}\label{eq:modelprior}
  \text{prior of }S=\frac{w(\abs{S})}{\bar w_K},\qquad w(k):=\frac{1}{2\binom{2p-1}{k}},\qquad
  \bar w_K:=\sum_{k=0}^{K}\binom{p}{k}w(k).
\end{equation}
This is the beta-binomial prior $\BB(1,p)$ \citep{ScottBerger2010} restricted to $\Om$: a probability $q$ is drawn from
$\BetaDist(1,p)$, and given $q$ each predictor enters independently with probability $q$. Its mean model size is about
one, and adding a predictor to a model of size $k$ multiplies the prior by $(k+1)/(2p-1-k)$, a cost of about
$\log\{2p/(k+1)\}$ in log prior odds. For example, with $p=3$ and $n\ge6$ the empty model has prior probability $\tfrac12$,
each one-predictor model $\tfrac1{10}$, and each larger model $\tfrac1{20}$.

\paragraph{Coefficients and noise variance.} Given $S$ and $\sigma^2$, the coefficients have the ridge-stabilised
$g$-prior \citep{Zellner1986} with $g=n$, and the variance an inverse-gamma prior:
\begin{equation}\label{eq:slab}
  \theta_S\mid S,\sigma^2\sim\Normal\big(0,\;n\sigma^2(G_S+I)^{-1}\big),\qquad \sigma^2\sim\IG(a_0,b_0),\quad a_0=b_0=0.01 .
\end{equation}
The factor $g=n$ makes the prior worth about one observation, and the unit ridge keeps it proper when $G_S$ is singular.

\subsection{The posterior, in closed form}\label{sec:posterior}

Integrating $\theta_S$ and $\sigma^2$ out gives the marginal likelihood of each model,
\begin{equation}\label{eq:marglik}
  m(S)=C_n\Big(\frac{\det H_S}{\det(G_S+I)}\Big)^{-1/2}(2b_0+\SSq_S)^{-\nu},\quad
  H_S:=(n+1)G_S+I,\ \ \SSq_S:=y^\top y-n\,z_S^\top H_S^{-1}z_S,
\end{equation}
with $\nu:=a_0+n/2$ and $C_n:=(2b_0)^{a_0}\Gamma(a_0+n/2)/\{\pi^{n/2}\Gamma(a_0)\}$ (here $\pi=3.14\ldots$). $\SSq_S$ is a
shrunken residual sum of squares; the determinant ratio is a penalty for size (an Occam factor), which in a
well-conditioned model multiplies $m(S)$ by about $(n+1)^{-1/2}$ for each extra predictor. The \emph{log score} of a model
and the \emph{posterior over models} are
\begin{equation}\label{eq:posterior}
  F(S):=\log\{w(\abs{S})/\bar w_K\}+\log m(S),\qquad \pis{S}=\frac{e^{F(S)}}{Z},\qquad Z:=\sum_{S'\in\Om}e^{F(S')} .
\end{equation}
The normalising constant $Z$ is the \emph{evidence}. Because $y$ is standardised, the posterior does not change if $\yr$
is shifted or rescaled. The \emph{inclusion probability} of predictor $j$ is $\PIP_j:=\sum_{S\ni j}\pis{S}$; the
\emph{highest posterior model} is $\HPM:=\argmax_S\pis{S}$; the \emph{median probability model} is
$\MPM:=\{j:\PIP_j>1/2\}$ \citep{BarbieriBerger2004}; and the \emph{posterior miss} is $1-\pis{\Sz}$.

\begin{lemma}[The posterior miss controls the summaries]\label{lem:summaries}
For every data set, $\max_{j\in\set{p}}\abs{\PIP_j-\ind{j\in\Sz}}\le1-\pis{\Sz}$. If $\pis{\Sz}>1/2$, then
$\HPM=\MPM=\Sz$.
\end{lemma}

The proof is two lines: every model that misses a true predictor, or contains a null one, differs from $\Sz$.

\subsection{Why sequential Monte Carlo, and why this start law}\label{sec:intuition}

The posterior is known up to one constant, so the difficulty is to find where its mass is in an enormous, rugged model
space. It is rugged for two reasons: correlated predictors create near-ties between models that swap a true predictor for
its proxy, and a predictor that is useless once the true ones are in can look useful on its own, because it \emph{leaks}
part of their signal. Sequential Monte Carlo starts many models (\emph{particles}) from an easy distribution close to the
target and moves them to the target in small steps, copying the particles the next distribution favours and spreading
them with Markov moves. The \emph{start law} is that easy distribution: it scores every subset of each block of correlated
predictors exactly, so it resolves a true predictor against its proxies, and it can be sampled exactly. What it gets wrong
is leakage across blocks. The local search from its mode undoes that, deleting the spurious predictors one per move; it
uses no random numbers, so it gives a second route to the true model that does not depend on how well the particles
mixed.

\begin{example}[Leakage, with six predictors]\label{ex:leak}
Take $n=100$, true model $\{1,3\}$ with both coefficients $1$, and unit noise. Predictor $2$ is a proxy of $1$
(correlation $0.93$); $4$ is moderately correlated with $3$ ($0.69$); $5$ is correlated with $1$ ($0.63$); and $6$ is a
proxy of $5$ ($0.89$). With blocks $\{1,2\}$, $\{3,4\}$, $\{5,6\}$, the five models with the largest posterior mass are:

\smallskip
\centerline{\small
\begin{tabular}{@{}lccccc@{}}
\toprule
model & $\{1,3,5\}$ & $\{1,3\}$ & $\{1,3,6\}$ & $\{1,3,4\}$ & $\{1,2,3\}$ \\
\midrule
start law $\pi_0$ & $\mathbf{0.517}$ & $0.198$ & $0.067$ & $0.012$ & $0.008$ \\
posterior $\pi(\cdot\mid y)$ & $0.042$ & $\mathbf{0.688}$ & $0.118$ & $0.051$ & $0.037$ \\
\bottomrule
\end{tabular}}
\smallskip

Within its own block, predictor $5$ explains much of $y$ because it carries part of predictor $1$'s signal, and the
start law cannot see that predictor $1$, in another block, already explains it: predictor $5$ leaks, and the start mode is
$\{1,3,5\}$. The proxy $2$, in predictor $1$'s own block, is handled correctly. The posterior puts $0.69$ on $\{1,3\}$, and
the local search from the start mode makes one move, deleting $5$, and stops at $\{1,3\}$, which has no improving neighbour. Result~3 proves this picture
in general: the start mode contains the true model and possibly some predictors that enter it wrongly (here
predictor~$5$), all members of a small \emph{leakage set} (Section~\ref{sec:reported}), and the search deletes them. (Computed with an independent re-implementation of the coded score,
start law and search on this data set; with six predictors the software itself would enumerate all $64$ models.)
\end{example}

\subsection{The algorithm, step by step}\label{sec:algorithm}

A run has seven stages (Figure~\ref{fig:pipeline}). A \emph{step} is one move of the sampler from one temperature to the
next.

\begin{figure}[ht]
\centering
\begin{tikzpicture}[node distance=5mm, every node/.style={font=\footnotesize},
  box/.style={draw, rounded corners, align=center, minimum height=10mm, text width=21mm},
  exact/.style={box, fill=blue!12, line width=0.9pt}, stoch/.style={box, fill=orange!18, dashed, line width=0.9pt},
  det/.style={box, fill=green!14, double}]
\node[exact] (pre) {1. standardise;\\ small $\Om$: exact};
\node[exact, right=of pre] (blk) {2. blocks of\\ $\le B$; tables};
\node[exact, right=of blk] (start) {3. start law $\pi_0$};
\node[stoch, right=of start] (cal) {4. calibrate,\\ then freeze};
\node[stoch, below=12mm of cal] (smc) {5. four islands\\ from $\pi_0$ to $\pi_1$};
\node[det, left=of smc] (ls) {6. two local\\ searches};
\node[det, left=of ls] (out) {7. report};
\draw[-{Stealth}] (pre)--(blk); \draw[-{Stealth}] (blk)--(start); \draw[-{Stealth}] (start)--(cal);
\draw[-{Stealth}] (cal)--(smc); \draw[-{Stealth}] (smc)--(ls); \draw[-{Stealth}] (ls)--(out);
\draw[-{Stealth}, dashed, thick] (start.south) -- node[right, pos=0.5]{start mode $\tSz$} (ls.north);
\draw[-{Stealth}, dotted, thick] (pre.south) |- node[left, pos=0.3]{exact case} (out.west);
\end{tikzpicture}
\caption{The stages of a run. Exact computations: solid boxes (1--3). Random: dashed boxes (4--5). Deterministic given
their inputs: double boxes (6--7). When $\Om$ is small the run goes from stage~1 straight to the report (dotted). The
second local search starts from the start mode and does not depend on the particles (dashed arrow).}
\label{fig:pipeline}
\end{figure}

\paragraph{Stage 1: inputs, the exact case and the check before the run.} The inputs are \emph{valid} when $n\ge3$, at
least one predictor is given and every entry is finite; the seed is an integer in $[1,2^{31}-1]$. Invalid inputs stop the
run. With $N_0:=\max(2500,n,p)$ and $\Nc_0:=\lceil512N_0/2500\rceil$, if $\abs{\Om}\le(5\Nc_0+4N_0)\,p$ the method scores every model
and computes everything exactly. Otherwise it estimates the memory the run will need and the worst-case time of the local search from the start
mode, and stops with status \texttt{resource\_limit} if the
estimate exceeds the allocation (8\,GiB and 24 hours by default).

\paragraph{Stage 2: blocks and block tables.} The method partitions the predictors into blocks $C_1,\dots,C_J$ of at most
$B:=\min(\lfloor\sqrt n\rfloor,\lfloor\sqrt p\rfloor,16)$ predictors, by complete-linkage clustering on $1-\abs{\hat r_{jl}}$,
cutting the tree top-down until every cluster fits and packing whole clusters into blocks in leaf order. The partition is a
deterministic function of $X$, and correlated predictors tend to share a block. Within each block it scores every subset
$U$ with $\abs{U}\le K$ and stores the \emph{block table} $h_k(U):=m(U)/m(\emptyset)$, with $h_k(\emptyset)=1$; a block has
at most $2^{16}$ subsets. It also builds a \emph{neighbour graph} of correlated predictors for the swap moves
(Appendix~\ref{app:params}).

\paragraph{Stage 3: the start law and the start mode.} The start law multiplies the size prior by the block tables,
\begin{equation}\label{eq:startlaw}
  \pi_0(S)=\frac{\gamma_0(S)}{\bar Z_0},\qquad \gamma_0(S):=\frac{w(\abs{S})}{\bar w_K}\prod_{k=1}^J h_k(S\cap C_k),\qquad
  \bar Z_0:=\sum_{S\in\Om}\gamma_0(S):
\end{equation}
the posterior one would have if the marginal likelihood, relative to the empty model, split into independent block
contributions. It sees every interaction within a block and none between blocks. A recursion over model sizes, combining
the blocks one at a time, computes $\bar Z_0$ exactly and gives an exact sampler for $\pi_0$ (draw the total size, then
the size of each block backwards, then a subset of each block). The \emph{start mode} $\tSz$ is a maximiser of $\pi_0$,
computed by the same recursion with maxima: the smallest size among the maximisers, then the lexicographically smallest
set. The posterior and the start law differ by the \emph{discrepancy}
$\Delta(S):=\log\{m(S)/m(\emptyset)\}-\sum_k\log h_k(S\cap C_k)$, the part of the log marginal likelihood that comes from
interactions between blocks: $\gamma_0(S)e^{\Delta(S)}=e^{F(S)}/m(\emptyset)$.

\paragraph{Stages 4 and 5: the sampler.} The sampler moves through the distributions
\begin{equation}\label{eq:path}
  \pi_\xi\ \propto\ \gamma_0\,e^{\xi\Delta},\qquad 0\le\xi\le1,
\end{equation}
from the start law ($\xi=0$) to the posterior ($\xi=1$), along a grid of temperatures $0=\grid0<\dots<\grid{\Ts}=1$ fixed
by the calibration. An island of $N$ particles starts from $N$ independent draws from $\pi_0$. At each step $t$ it
\begin{enumerate}[label=(\alph*),leftmargin=2em,itemsep=0pt,topsep=2pt]
  \item \emph{reweights} particle $v$ by $\omega^{(t)}_v:=\exp\{(\grid t-\grid{t-1})\Delta(S_v)\}$;
  \item \emph{resamples} the particles multinomially, in proportion to these weights;
  \item \emph{moves} every particle by one round of Markov moves that leave $\pi_{\grid t}$ invariant.
\end{enumerate}
A round is an inclusion update (Gibbs) for each predictor, one uniform swap, $B$ swaps along the neighbour graph, and one
independence move that proposes a fresh draw from $\pi_0$. At $\xi=1$ it runs $\Lbi:=\max(4,\lceil\log_2(n+1)\rceil)$ burn-in rounds
and four measured rounds. The island's inclusion-probability estimate $\PIPhat_{r,j}$ averages, over particles and measured
rounds, the conditional inclusion probability of the Gibbs update rather than the 0/1 indicator; its evidence estimate is
$\Zhat_r:=m(\emptyset)\,\bar Z_0\prod_{t=1}^{\Ts}\{N^{-1}\sum_v\omega^{(t)}_v\}$, the exact normaliser of the start law times
the product of the average reweighting factors.

The \emph{calibration} (stage 4) chooses the grid and the effort before the final run, from one pilot island and a few
rounds of four short diagnostic islands; the grid refines the pilot's temperatures so that each step keeps an
effective-sample-size fraction, computed on a set of candidate models, of at least $0.90$. The calibration starts at
tier $i=0$ and moves at most once to tier $i=1$, which multiplies all particle numbers by four, when one of its
diagnostics fires (Appendix~\ref{sec:supp-cal} gives the full procedure). Three facts are all that the results use: the calibration's outcome is a function of the data and of its own
random numbers; it is \emph{frozen} before the final run draws any random number; and it \emph{succeeds} when it freezes a
grid, by its main construction or a fallback (otherwise the run stops without estimates). The \emph{final run} (stage 5)
consists of four independent \emph{islands} of $N_i=4^iN_0$ particles at the frozen tier $i$, on the frozen grid, each with
its own random stream.

\paragraph{Stage 6: the local searches.} Two models are \emph{neighbours} if one comes from the other by deleting a
predictor, adding one (while the size is below $K$), or swapping a member for a non-member. A local search scans every
neighbour of the current model and moves to the best one (ties broken lexicographically) if its log score is strictly
higher, for at most a capped number of moves. It stops at the first scan with no strict improvement and \emph{certifies}
its end point; after its last allowed move it scans once more, and if a neighbour still improves it stops there,
uncertified. Two searches run after the islands: from the best model found so far, with move cap
$\ell_n:=\lceil\log_2(n+p+1)\rceil$; and from the start mode $\tSz$, with move cap $\ell_n+\abs{\tSz}$.

\subsection{What a run reports}\label{sec:outputs-method}

A run either \emph{produces output}, with status ``exact'' (enumeration) or ``completed'' (sampler; a \emph{completed
run}), or stops without estimates. A run stopped by its time guard reports status \texttt{resource\_limit} and returns no
estimates; anything else such a run returns is flagged incomplete, and the results do not cover it. (A constant
response, or no retained predictor, stops the run with the empty model; invalid inputs stop it with no model.) The method records every model
it scores; the \emph{reported model} $\Shat$ is the best of them by log score (exact ties to the smaller model, then
lexicographically). It is \emph{certified} when the method enumerated $\Om$, or when $\Shat$ ends a local search whose last
scan found no strict improvement. An exact run reports $\Shat$, every $\PIP_j$, the median probability model and $\log Z$.
A completed run pools its islands with equal weights,
\begin{equation}\label{eq:pool}
  \PIPhat_j:=\frac14\sum_{r=1}^4\PIPhat_{r,j},\qquad \Zhat:=\frac14\sum_{r=1}^4\Zhat_r,
\end{equation}
and reports $\Shat$, the $\PIPhat_j$, the estimated median probability model $\MPMhat:=\{j:\PIPhat_j>1/2\}$, $\log\Zhat$
with a one-sided $95\%$ lower bound $\log(0.05\,\Zhat)$ for $\log Z$ (Appendix~\ref{app:params}), a flag saying whether
$\Shat$ is certified, and a list $\mathcal T$ of at most $40p+1$ models with their log scores, the best among those it scored. Islands are not
weighted by their evidence estimates. Diagnostics, which never change an output, are listed in
Appendix~\ref{app:params}.

\subsection{The constants}\label{sec:constants}

Table~\ref{tab:params} lists the constants of the method that the results use or that a reader needs. Appendix~\ref{app:params}
gives all of them, with their code locations.

\begin{table}[ht]
\centering\small
\setlength{\tabcolsep}{4pt}
\caption{The main constants of the method. Every one is a numerical constant or a function of $(n,p)$; the example
column uses $n=200$, $p=10^4$.}\label{tab:params}
\begin{tabular}{@{}>{\raggedright\arraybackslash}p{0.34\linewidth}l>{\raggedright\arraybackslash}p{0.3\linewidth}r@{}}
\toprule
Constant & Symbol & Value & Example \\
\midrule
size cap & $K$ & $\min(p,\lfloor n/2\rfloor)$ & 100 \\
model prior & $w(k)$ & $\tfrac12/\binom{2p-1}{k}$, renormalised to $\Om$ & \\
coefficient prior scale; ridge & $g$ & $n$; unit ridge & 200 \\
inverse-gamma shape and scale & $a_0,b_0$ & $0.01$, $0.01$ & \\
block size cap & $B$ & $\min(\lfloor\sqrt n\rfloor,\lfloor\sqrt p\rfloor,16)$ & 14 \\
particles per island (tier $i$) & $N_i$ & $4^i\max(2500,n,p)$ & $10^4$; $4\cdot10^4$ \\
islands; measured rounds & & $4$; $4$ & \\
burn-in rounds at $\xi=1$ & $\Lbi$ & $\max(4,\lceil\log_2(n+1)\rceil)$ & 8 \\
moves per round & & $p$ inclusion updates, $1$ uniform swap, $B$ neighbour swaps, $1$ independence move & 10,016 \\
calibration grid target & & restricted effective-sample-size fraction $\ge0.90$ & \\
exact-case threshold & & enumerate if $\abs{\Om}\le(5\Nc_0+4N_0)\,p$ & $5.0\times10^8$ \\
base move cap & $\ell_n$ & $\lceil\log_2(n+p+1)\rceil$ & 14 \\
move cap from the start mode & & $\ell_n+\abs{\tSz}$ & \\
returned list & $\mathcal T$ & best $40p+1$ models by $F$ & 400,001 \\
pooling of islands & & equal weights & \\
\bottomrule
\end{tabular}
\end{table}

\subsection{What the results use from the method}\label{sec:conventions}

The results of Section~\ref{sec:results} use five objects of a run: the start law and its mode (stage 3), the reported
model and its certification (stages 6--7), the pooled estimates~\eqref{eq:pool}, the returned list $\mathcal T$, and the
event $\Acal$ that the method uses the sampler and its calibration succeeds. A run in $\Acal$ may still stop later, at
its time guard; every completed run lies in $\Acal$. They also
use three conventions, which fix what ``the method'' means and idealise the computer; they are not statistical
assumptions, and Appendix~\ref{app:params} lists them in full. Exact arithmetic and ideal random numbers are also
presupposed by Song and Liang's statements about their own sampler (their supplement, pp.~9--10, 12--13).
\begin{itemize}[leftmargin=1.5em,itemsep=1pt]
  \item \phantomsection\label{asm:C1}\emph{Convention \as{C1} (the method):} exact arithmetic, the default configuration,
    valid inputs and seed. A run that meets it is \emph{admissible}; a sample size $n$ is admissible if some data set of
    that size is valid. Statements hold for every valid seed and every outcome of the calibration.
  \item \phantomsection\label{asm:C2p}\emph{Convention \as{C2p} (delivery):} an idealised
    machine that never overflows or fails, ideal random numbers, an allocation of at least the memory and time that the method's own polynomial
    requirement formulas require, and the code's two fixed integer limits on the move caps lifted (see the remark after
    Theorem~\ref{thm:reported}).
  \item \phantomsection\label{asm:C3}\emph{Convention \as{C3} (sampler):} exact arithmetic, ideal random numbers
    independent of the data and of the calibration, an exact start sampler, and Markov moves that do what their definitions
    say.
\end{itemize}
